%=============================================================================
% Appendices
%=============================================================================

\appendix

\section{Capability Interaction Patterns}

\subsection{Simple Capability Invocation}

\begin{enumerate}[label=\arabic*.]
\item Agent Decision: ``Search for AI architecture papers''
\item Contact Layer: Route to \texttt{search\_web} capability
\item Provider Selection: Choose Provider A (fast)
\item Execution Runtime: Invoke Provider A API
\item Result: Return search results to agent
\end{enumerate}

\subsection{Multi-Capability Workflow}

\begin{enumerate}[label=\arabic*.]
\item Agent Decision: ``Write a blog post about AI''
\item Contact Layer: Decompose into capability sequence:
  \begin{enumerate}[label=\alph*.]
  \item \texttt{search\_web}(``AI topics'')
  \item \texttt{summarize}(search\_results)
  \item \texttt{generate\_text}(blog\_outline)
  \item \texttt{edit\_text}(generated\_draft)
  \end{enumerate}
\item Execution: Execute each capability in sequence
\item Coordination: Pass results between capabilities
\item Result: Return final blog post
\end{enumerate}

\subsection{Error Handling Pattern}

\begin{enumerate}[label=\arabic*.]
\item Agent Request: invoke capability X
\item Contact Layer: Route to Provider A
\item Execution: Provider A fails (timeout)
\item Contact Layer: Retry with Provider B
\item Execution: Provider B succeeds
\item Result: Return result from Provider B
\end{enumerate}

\section{Capability Definition Schema}

\subsection{Capability Schema}

\begin{lstlisting}[style=yamlstyle]
capability_schema:
  name: string
  version: string
  description: string
  input_schema:
    type: object
    properties: {}
    required: []
  output_schema:
    type: object
    properties: {}
  execution:
    endpoint: string
    method: enum [GET, POST, PUT, DELETE]
    timeout_ms: integer
    retry_policy:
      max_retries: integer
      backoff_ms: integer
  provider:
    name: string
    version: string
  metadata:
    category: enum [AI, API, Tool, Device]
    tags: [string]
    cost_per_call: number
\end{lstlisting}

\subsection{Example Capability}

\begin{lstlisting}[style=yamlstyle]
name: control_device
version: "1.0"
description: Control an IoT device
input_schema:
  type: object
  properties:
    device_id:
      type: string
    action:
      type: string
      enum: [turn_on, turn_off, set_level]
    level:
      type: number
      description: Level for set_level action
  required: [device_id, action]
output_schema:
  type: object
  properties:
    status:
      type: string
      enum: [success, failure]
    device_state:
      type: object
    timestamp:
      type: string
execution:
  endpoint: /v1/devices/{device_id}/control
  method: POST
  timeout_ms: 5000
  retry_policy:
    max_retries: 2
    backoff_ms: 1000
provider:
  name: iot_controller
  version: "2.1"
metadata:
  category: Device
  tags: [iot, physical, automation]
  cost_per_call: 0.001
\end{lstlisting}

\section{Contact Layer Interface Specification}

\subsection{Routing Interface}

\begin{lstlisting}[style=typescriptstyle]
interface ContactLayer {
  /**
   * Route a capability request to an appropriate provider.
   */
  route(request: CapabilityRequest): Promise<RoutingDecision>;

  /**
   * Evaluate policies for capability invocation.
   */
  evaluatePolicy(
    capability: string,
    context: ExecutionContext
  ): Promise<PolicyDecision>;

  /**
   * Select provider based on context.
   */
  selectProvider(
    capability: string,
    context: ExecutionContext
  ): Promise<Provider>;

  /**
   * Choose execution environment.
   */
  selectEnvironment(
    capability: string,
    provider: Provider
  ): Promise<Environment>;
}

interface RoutingDecision {
  capability: string;
  provider: Provider;
  environment: Environment;
  policyCompliant: boolean;
  estimatedLatency: number;
  estimatedCost: number;
}
\end{lstlisting}

\subsection{Capability Request Structure}

\begin{lstlisting}[style=typescriptstyle]
interface CapabilityRequest {
  // Capability identification
  capability: string;
  version?: string;

  // Input data
  input: Record<string, any>;

  // Execution context
  context: {
    sessionId: string;
    requestId: string;
    userId?: string;
    priority?: 'low' | 'medium' | 'high';
    deadline?: Date;
  };

  // Routing preferences
  preferences?: {
    preferredProvider?: string;
    maxLatency?: number;
    maxCost?: number;
    requirePrivacy?: boolean;
  };
}
\end{lstlisting}
